\section{Conclusion}

This paper has studied international standardization in a three-country model with network externalities, coalition formation, and asymmetric participation capacity. The key distinction developed here is between \emph{formal accession} to a common standard and \emph{effective participation} in that standard. Mutual recognition may eliminate the baseline technology-gap cost within a recognition area, yet countries can still differ in their ability to comply with, test, certify, and implement the common standard. Once that distinction is introduced, the equilibrium logic of international standardization changes in economically meaningful ways.

The analysis delivers three main conclusions. First, full international standardization may be globally efficient without arising in equilibrium. A low-capacity outsider may add too little to the common installed base, relative to the additional competitive pressure its accession creates in incumbent-member markets, for full harmonization to be privately attractive to the member governments. Second, stronger network effects do not necessarily favor full harmonization uniformly. On a verified parameter region with pronounced participation-capacity asymmetry, stronger network effects can instead stabilize an exclusionary two-country standards bloc. Third, once participation capacity is endogenized, capacity-building investment and conditional transfer policy can, on the verified compact domain used in Section~5, restore the missing accession margin more effectively than policies that merely preserve fragmented national standards.

These results contribute to the literature in a simple but useful way. Relative to the classic compatibility literature, the paper shows that compatibility and participation should not be treated as synonymous. Relative to the international standardization literature, it identifies a mechanism through which inefficient exclusion can persist even when mutual recognition is technologically feasible. Relative to recent empirical work on standards, trade, and innovation, it offers a tractable framework for thinking about how asymmetries in compliance and implementation capacity may shape the formation and persistence of standards blocs.

The policy message is equally direct, but it should be read with the qualification emphasized in Section~5. If fragmentation persists because weaker countries cannot participate effectively in a harmonized regime, then the relevant intervention is not limited to the legal act of recognition itself. What matters is whether participation capacity can be raised through investment in testing, certification, conformity assessment, and related implementation capabilities. In the model, such policies relax the actual bottleneck in the regime-selection game, whereas policies that merely preserve fragmented national standards leave that bottleneck in place. This suggests that capacity-oriented assistance may be a more effective response to inefficient fragmentation than policies centered solely on preserving domestic standard autonomy.
This policy conclusion should be interpreted in the same qualified sense as
the threshold analysis in Section~5: it is established on the verified compact
domain reported in Appendix~B and under the corresponding crossing
conditions, rather than as an unrestricted claim over all admissible parameter
values.

The paper also points to several natural extensions. One would be to endogenize coalition formation more fully, allowing for simultaneous three-country negotiation or asymmetric bilateral bargaining involving country $C$ rather than the benchmark accession game studied here. Another would be to allow firms to choose not only quantities but also whether to adopt foreign or international standards strategically. A third would be to introduce partial interoperability, rather than the sharp distinction between exclusion and full membership used here. A fourth extension would be to connect the theoretical threshold conditions more directly to empirical measures of participation capacity, for example using data on standardization activity, certification infrastructure, or international standard adoption. Such extensions would help clarify how formal standardization institutions interact with implementation capacity in practice.

More broadly, the paper suggests that the economics of international standardization should pay greater attention to the difference between joining a rule system and being able to operate effectively within it. In markets shaped by network effects and compatibility decisions, that distinction can determine whether harmonization becomes genuinely inclusive or remains organized around exclusionary blocs.
